% Author: Writer agent · 2026-11-14
% Task W2: Chapter 6 (generic pipeline). Sources: research-notes.md, Ch. 6 (+ Ch. 2, 3 keys for blocks 1 and 3).
% The eight-block division is the BOOK'S OWN educational simplification (see Gaps in research-notes.md).

\chapter{שמונה תחנות בדרך לשיר הבא}\label{ch:6}

\begin{scene}
21:00. נועה מבשלת שקשוקה, ידיים מלאות עגבניות. היא לוחצת על הטלפון עם המרפק: „תני משהו". חצי שנייה אחר כך כבר מתנגן שיר, והוא מתאים בדיוק לשקשוקה.

„רגע", היא אומרת למחבת, „מה קרה בחצי השנייה הזאת?"

המחבת, כמו הטלפון בפרק~\ref{ch:1}, לא עונה. אבל הפרק הזה כן.
\end{scene}


\section{מפה לימודית, לא שרטוט של מפעל}

איור~\ref{fig:pipeline} מחבר את כל מה שלמדנו לתהליך אחד. חשוב לומר ביושר: החלוקה לשמונה תחנות היא פישוט לימודי שלנו, ולא תיאור של המערכת של חברה כלשהי. המקורות תומכים בדפוס הכללי. לפי המדריך של \en{Google}, בדפוס נפוץ המערכת מתחילה מאוסף ענק ומצמצמת אותו לקבוצה קטנה בהרבה של מועמדים, אחר כך מדרגת אותם, ולבסוף מתחשבת באילוצים נוספים לפני הדירוג הסופי \cite{google2026overview}. מאמר של \en{YouTube} מ-2016 מתאר מבנה דומה של שני שלבים: מודל ליצירת מועמדים ומודל נפרד לדירוג \cite{covington2016deep}. גם לולאת המשוב מתועדת במחקר \cite{chaney2018algorithmic}. שלוש התחנות הראשונות הן מסגור סביר שלנו, אבל אין מקור אחד שמונה את כל שמונה התחנות.

\begin{herefigure} % placement via the class (Prompt 18; audit A2-13)
  \centering
  % Eight-block recommendation process: EDUCATIONAL SIMPLIFICATION (Boss, 2026-11-14).
% Not the internal architecture of any platform. Editable vector TikZ.
% Layout: snake, right-to-left. Top row 1→4 (right to left), down to 5, bottom row 5→8
% (left to right), then the feedback loop 8→1 on the right edge.
% RTL rules (PRD §6): Hebrew labels coloured via \hebcolor (never [heb] + colour);
% number badges are digits only and need no RTL wrapping.
\begin{tikzpicture}[
    x=1cm,y=1cm,
    % Prompt 13 palette: group fill + border; charcoal labels on every fill (contrast >= 5.4:1)
    blk/.style={draw=#1,line width=1pt,fill=#1!28!white,rounded corners=2.5mm,
      minimum width=3.1cm,minimum height=1.45cm,align=center,
      font=\headfont\small,inner sep=2pt},
    num/.style={circle,fill=#1,minimum size=5.2mm,inner sep=0pt,
      font=\sffamily\bfseries\scriptsize,text=charcoal},
    flow/.style={-{Stealth[length=2.6mm,width=2.2mm]},line width=1pt,draw=charcoal!80},
  ]
  % --- blocks (centre x: 13.55, 9.55, 5.55, 1.55; rows y = 0 and -2.75; total width ≈ 16 cm)
  \node[blk=brandgreen] (b1) at (13.55,0)    {\hebcolor{charcoal}{איסוף}\\\hebcolor{charcoal}{אותות}};
  \node[blk=brandgreen] (b2) at (9.55,0)    {\hebcolor{charcoal}{ניקוי ואיחוד}\\\hebcolor{charcoal}{נתונים}};
  \node[blk=brandgreen] (b3) at (5.55,0)     {\hebcolor{charcoal}{ייצוג משתמשים}\\\hebcolor{charcoal}{ושירים}};
  \node[blk=charcoal] (b4) at (1.55,0)    {\hebcolor{charcoal}{יצירת}\\\hebcolor{charcoal}{מועמדים}};
  \node[blk=charcoal] (b5) at (1.55,-2.75){\hebcolor{charcoal}{דירוג}};
  \node[blk=charcoal] (b6) at (5.55,-2.75) {\hebcolor{charcoal}{סינון}\\\hebcolor{charcoal}{וגיוון}};
  \node[blk=brandpink](b7) at (9.55,-2.75){\hebcolor{charcoal}{הצגת}\\\hebcolor{charcoal}{המלצות}};
  \node[blk=brandpink](b8) at (13.55,-2.75){\hebcolor{charcoal}{משוב}\\\hebcolor{charcoal}{ולמידה}};
  % --- number badges on each block's top-right corner (reading start in RTL)
  % digit colour per badge: charcoal on green/pink, light on charcoal (digits only, so text= is safe)
  \foreach \b/\n/\c/\t in {b1/1/brandgreen/charcoal,b2/2/brandgreen/charcoal,b3/3/brandgreen/charcoal,b4/4/charcoal/pagebg,b5/5/charcoal/pagebg,b6/6/charcoal/pagebg,b7/7/brandpink/charcoal,b8/8/brandpink/charcoal}
    \node[num=\c,text=\t] at (\b.north east) {\n};
  % --- flow
  \draw[flow] (b1) -- (b2);
  \draw[flow] (b2) -- (b3);
  \draw[flow] (b3) -- (b4);
  \draw[flow] (b4) -- (b5);
  \draw[flow] (b5) -- (b6);
  \draw[flow] (b6) -- (b7);
  \draw[flow] (b7) -- (b8);
  % --- feedback loop 8 → 1 around the right edge
  \draw[flow,draw=brandpink,dashed,line width=1.3pt] (b8.east) -- ++(0.45,0) |- (b1.east);
  % Prompt 24 (users' decision on QA finding D1-04): label moved 5.5 pt right → 2.0 mm gap to the dashed line.
  % The phantom copy at the old place keeps the picture's bounding box identical (blocks do not move);
  % the visible label is an overlay, so it does not widen the picture (no overfull line).
  \node[anchor=center,font=\headfont\footnotesize,rotate=90] at (15.7,-1.375) {\phantom{לולאת משוב}};
  \node[overlay,anchor=center,font=\headfont\footnotesize,rotate=90] at (15.894,-1.375)
    {\hebcolor{charcoal}{לולאת משוב}}; % small text: charcoal, not pink
  % --- legend: three groups (no digit ranges: "1–3" displays as "3–1" in RTL)
  \begin{scope}[shift={(0,-4.35)}]
    \foreach \x/\c/\t in {13.55/brandgreen/{נתונים},9.55/charcoal/{מודל},5.55/brandpink/{משתמש}}{
      \fill[\c] (\x+1.05,0) circle (1.6mm);
      \node[anchor=east,font=\headfont\footnotesize] at (\x+0.85,0) {\hebcolor{charcoal}{\t}};}
  \end{scope}
\end{tikzpicture}

  \caption{תהליך ההמלצה בשמונה תחנות: פישוט לימודי של הספר, ולא הארכיטקטורה הפנימית של פלטפורמה כלשהי \educational}
  \label{fig:pipeline}
\end{herefigure}

\section{התחנות, בערב של נועה}

\begin{enumerate}
  \item \textbf{איסוף אותות.} כל פעולה של נועה נרשמת: הדילוג של 23:00 אתמול מפרק~\ref{ch:5}, ההאזנה לשיר האינדי מפרק~\ref{ch:4}. כפי שראינו בפרק~\ref{ch:2}, לפי \en{Spotify} פעולות כאלה משפיעות על האופן שבו היא מפרשת את הטעם \cite{spotify2026understanding}.
  \item \textbf{ניקוי ואיחוד נתונים.} מאחדים את האירועים מהטלפון ומהמחשב, ומנסים לזהות רעש. למשל, שיר שהתנגן שעה כשהטלפון נשכח בכיס הוא לא בהכרח שיר אהוב.
  \item \textbf{ייצוג משתמשים ושירים.} נועה וכל שיר הופכים לייצוג וקטורי, כלומר לנקודה במרחב משותף, כמו בפרק~\ref{ch:3} \cite{google2026candidate}.
  \item \textbf{\term{יצירת מועמדים}{candidate generation}.} מתוך הקטלוג העצום נבחרת קבוצה קטנה בהרבה \cite{google2026overview}: שירים שיוני אוהב, שירים שדומים לשיר האינדי. (בנפרד, לפי \en{Spotify}, האלגוריתמים שלה מתחשבים גם במגמות ובטעם הכללי של המשתמשים, בלי לפרט באיזה שלב \cite{spotify2026understanding}.)
  \item \textbf{\term{דירוג}{ranking}.} מודל נפרד מעריך כל מועמד ומסדר את הרשימה \cite{covington2016deep}. בפישוט שלנו, השאלה היא בערך: „מה הסיכוי שנועה תשמע את זה עד הסוף, עכשיו?"
  \item \textbf{סינון וגיוון.} ב\term{דירוג מחדש}{re-ranking} אפשר להסיר פריטים שהמשתמש סימן במפורש שאינו אוהב, לתת דחיפה לתוכן טרי, ולעזור להבטיח גיוון, רעננות והוגנות \cite{google2026overview}. כך אפשר למנוע מצב שבו יתנגנו עשרה שירים של אותה זמרת ברצף. לפי הפישוט שלנו, זה גם מקום טבעי לחקירה מפרק~\ref{ch:5}.
  \item \textbf{הצגת המלצות.} השיר הראשון מתנגן ליד השקשוקה, ושאר הרשימה מחכה בתור.
  \item \textbf{משוב ולמידה.} התגובה של נועה הופכת לאות חדש. לפי \en{Spotify}, לחיצה על „לא מעוניין", אגודל למטה או הסתרה מובילה לפחות המלצות דומות \cite{spotify2026understanding}.
\end{enumerate}

החץ המקווקו בין התחנה האחרונה לראשונה הוא \term{לולאת משוב}{feedback loop}. מערכות המלצה מאומנות לעיתים קרובות על נתונים של משתמשים שכבר ראו המלצות. לפי צ'ייני ועמיתיו, זה יוצר לולאת משוב מזיקה, ובסימולציות שלהם היא הפכה את התנהגות המשתמשים לאחידה יותר בלי להגדיל את התועלת שלהם \cite{chaney2018algorithmic}. במה זה נוגע לנועה? על כך בפרק הבא.

\begin{inshort}
\begin{itemize}
  \item שמונה התחנות הן פישוט לימודי של הספר, ולא המערכת של חברה כלשהי.
  \item הליבה מתועדת במקורות: יצירת מועמדים, דירוג, ואז סינון וגיוון.
  \item כל תגובה של נועה חוזרת לתחנה הראשונה וסוגרת לולאת משוב.
\end{itemize}
\end{inshort}
